\chapter{Introduction}
\label{chap:intro}

\subsubsection{Background and Motivation}

In practice, quantum key distribution is severely limited by distance because signals inevitably degrade due to noise in the communication channel (like optical fibers or free space). To overcome this, standard quantum repeaters are used to establish long-distance secure links by performing entanglement distillation (extracting pure, perfect entangled states from many noisy ones) and reliably swapping them across multiple intermediate nodes. However, there are bound entangled states from which no pure entanglement can be distilled, but from which a secure cryptographic key can still be extracted \cite{horodecki2005secure}. This raised an important question: Could a specialized quantum key repeater be built to bypass the need for pure entanglement distillation, allowing one to swap and extend the secret keys from these highly noisy states over long distances? \cite{limitations2015} 

Quantum channels provide the fundamental mathematical framework for describing the evolution of quantum systems and the transmission of quantum information. One particular class of noise is modeled by completely copositive (PPT) channels, which, when applied to a subsystem, ensure that the global state has a positive semidefinite partial transpose. Consider a quantum key repeater scenario where the point-to-point communication links are modeled by such PPT channels. A minimal requirement for successful key distribution in this scenario is the preservation of entanglement across the shared bipartite states. Interestingly, in \cite{christandl2012ppt} the PPT$^2$ conjecture was formulated, stating that the composition of any such PPT channels is an entanglement-breaking channel, which completely annihilates the desired quantum correlations when applied to a subsystem. If true, this implies that a signal traversing two such sequential noise channels would have its quantum correlations completely annihilated, rendering the resulting end-to-end state fully separable.

Beyond its immediate practical implications for quantum network architectures, determining the exact number of iterations required for a quantum channel to become entanglement-breaking is also of mathematical interest. Significant research has recently been dedicated to characterizing the behavior of the entanglement-breaking property under compositions. Notably, the PPT$^2$ conjecture was proven to be true for qutrit systems in \cite{christandlWhenComposedMaps2019}, its general validity for arbitrary finite dimensions remains to be an open problem in quantum information theory. 

\subsubsection{Outline}
To establish the necessary framework for analyzing these phenomena, Chapter 2 provides a brief introduction to the core concepts of quantum information theory, alongside the basic mathematical foundations required. Chapter 3 focuses on the concepts of entanglement and positive maps. Chapter 4 subsequently offers a more detailed overview of quantum channels and some subclasses with a focus on their matrix representations. Chapter 5 introduces the core notions of eventually entanglement-breaking channels and entanglement-breaking indices. The chapter reviews some of the existing literature and highlights some recent theoretical developments in the field. Finally, Chapter 6 establishes some necessary and separately sufficient analytical conditions for determining whether a given channel is entanglement-breaking.
